% Generated draft: requires mathematical and source review.
\documentclass[11pt]{article}
\usepackage[margin=25mm]{geometry}
\usepackage{fontspec}
\setmainfont{Latin Modern Roman}
\usepackage{amsmath,amssymb,mathtools,unicode-math}
\setmathfont{Latin Modern Math}
\usepackage{enumitem,needspace,xurl,hyperref}
\setlength{\parindent}{0pt}
\setlength{\parskip}{5pt}
\setlength{\emergencystretch}{3em}
\newcommand{\N}{\mathbb{N}}
\newcommand{\Z}{\mathbb{Z}}
\newcommand{\Q}{\mathbb{Q}}
\newcommand{\R}{\mathbb{R}}
\newcommand{\C}{\mathbb{C}}
\newcommand{\F}{\mathbb{F}}
\newcommand{\A}{\mathbb{A}}
\newcommand{\PP}{\mathbb{P}}
\newcommand{\E}{\mathbb{E}}
\newcommand{\eps}{\varepsilon}
\newcommand{\dd}{\,\mathrm{d}}
\begin{document}
\section*{UnsolvedMath ID: 801011}
\textbf{Source problem code:} catalogue entry
\subsection{Definitions and mathematical statement}
The supplied record concerns a map \(A\mapsto I(A)\), but supplies neither its definition nor its domain or codomain. It also uses \(H^f\) without defining the admissible functions \(f\) or the associated set function. Consequently, the following four questions retain unresolved notation and require the missing source context before they constitute fully specified mathematical problems.

\begin{enumerate}
\item Determine whether
\[
\bigl|\operatorname{ran}(I)\bigr|=\bigl|\mathcal P(\R)\bigr|,
\]
where \(\mathcal P(\R)\) is the power set of \(\R\), interpreting the source's \(2^{\mathbb R}\) in its usual power-set sense. Here \(\operatorname{ran}(I)=\{I(A):A\in\operatorname{dom}(I)\}\), with \(\operatorname{dom}(I)\) unspecified. The source adds ``Seems yes under GCH''; this is a supplied tentative qualification, not an established implication asserted here. GCH denotes the generalized continuum hypothesis.

\item The source states that for a closed set \(C\), \(I(C)\) is arithmetic, of complexity \(\mathbf{\Sigma}^{0}_{2}\). It asks for the complexity of
\[
\mathcal S=\bigl\{\varphi:\varphi\text{ is a }\mathbf{\Sigma}^{0}_{2}\text{ formula and defines }I(C)\text{ for some closed }C\bigr\},
\]
and specifically asks whether \(\mathcal S\) is \(\mathbf{\Sigma}^{1}_{2}\)-complete. The source further states that \(\mathcal S\) is \(\mathbf{\Pi}^{1}_{1}\)-hard. These complexity assertions are supplied claims. The ambient space in which \(C\) is closed, the objects defined by the formulas, permitted parameters, coding of formulas, and reductions used for hardness and completeness are not specified. Thus \(\mathcal S\) is presently a schematic description, not a fixed coded set.

\item Determine whether the assertion described in the source as ``Borel hierarchy is proper w.r.t. the range of \(I\)'' is independent of ZFC, the Zermelo--Fraenkel axioms with choice. The source does not specify the underlying topological space, how the range of \(I\) is compared with Borel classes, or what precise strictness assertion is intended. A definite sentence whose independence is being asked about therefore remains to be supplied.

\item Determine whether it is consistent with ZFC that, for every admissible \(f\) and every \(A\subseteq\R\), positivity of \(H^f(A)\) implies the existence of a subset \(A_0\) with finite positive \(H^f\)-value. If ``a subset \(A_0\)'' means a subset of \(A\), the intended assertion has the conditional form
\[
\forall f\;\forall A\subseteq\R\;\Bigl(H^f(A)>0\ \Longrightarrow\ \exists A_0\subseteq A\;\bigl(0<H^f(A_0)<\infty\bigr)\Bigr).
\]
The quantifier over \(f\) must be restricted to the intended, presently unspecified class. The displayed formulation explicitly records the natural interpretation \(A_0\subseteq A\); the supplied wording does not state that inclusion explicitly. No definition of \(H^f\), including any conventions governing its values, is supplied.
\end{enumerate}
\subsection{Short English statement}
Problem 23 asks four questions about an unspecified assignment \(A\mapsto I(A)\): the size of its range, the definability complexity of outputs arising from closed sets, independence of a Borel-hierarchy strictness assertion, and consistency of a finite-positive-subset property for \(H^f\). Missing definitions and coding conventions prevent a complete formal statement. The registry labels the problem open; that label and the accompanying complexity and GCH comments are supplied information requiring review.
\subsection{Sources}
\begin{itemize}
\item Local UnsolvedMath registry, supplied complete source record: permanent ID \(801011\), exact source problem code \texttt{catalogue entry}, title \emph{Open Problems in Mathematical Logic --- Problem 23}. The supplied proposer attribution is Barmpalias, Gao et al., with proposed year 2026. The registry status is \texttt{open}; it is not an independently verified literature assessment.
\item Supplied source attribution: Barmpalias, Gao et al., \emph{Open Problems in Mathematical Logic} (2026), Problem 23. Supplied source URL: \url{https://arxiv.org/abs/2608.26628}. The record reports access on 2026-10-16 and extraction by \texttt{arxiv\_latex}. These are supplied provenance details; the source has not been independently visited or verified for this draft. The supplied excerpt omits the definitions needed to resolve the notation.
\end{itemize}
\end{document}
